\chapter*{Заключение}                       % Заголовок
\addcontentsline{toc}{chapter}{Заключение}  % Добавляем его в оглавление

%% Согласно ГОСТ Р 7.0.11-2011:
%% 5.3.3 В заключении диссертации излагают итоги выполненного исследования, рекомендации, перспективы дальнейшей разработки темы.
%% 9.2.3 В заключении автореферата диссертации излагают итоги данного исследования, рекомендации и перспективы дальнейшей разработки темы.
%% Поэтому имеет смысл сделать эту часть общей и загрузить из одного файла в автореферат и в диссертацию:

В настоящей работе метод молекулярной динамики был применён для исследования фазовых переходов и сопутствующих процессов. В ходе расчётов получены зависимости плотности и коэффициента диффузии от давления для 2,2,4-триметилпентана ($C_8H_{18}$) и 1,1-дифенилэтана ($C_{14}H_{14}$). Для двухфазной системы «метан–этан» определён парахор и рассчитаны изотермы адсорбции. Для наночастиц железа в диапазоне температур кластера 200–2500 К вычислены коэффициенты термической аккомодации, теплоотдачи и прилипания.

По результатам работы сформулированы следующие выводы:

%% Согласно ГОСТ Р 7.0.11-2011:
%% 5.3.3 В заключении диссертации излагают итоги выполненного исследования, рекомендации, перспективы дальнейшей разработки темы.
%% 9.2.3 В заключении автореферата диссертации излагают итоги данного исследования, рекомендации и перспективы дальнейшей разработки темы.

\begin{enumerate}
  \item Установлено, что хотя прямые расчёты плотности и коэффициента самодиффузии демонстрируют расхождение с экспериментальными данными, последующая коррекция расчётных плотностей позволяет добиться более близкого к референтным значениям согласия по коэффициенту диффузии.
  \item Воспроизведено значение парахора чистного этана в диапазоне давлений 4–40 атм, что подтверждает возможность использования метода молекулярной динамики для прогнозирования экспериментальных значений поверхностного натяжения.
  \item Проведена верификация формулы Гиббса для оценки адсорбции на плоской поверхности в двухфазной двухкомпонентной системе. Формула, использующая только экспериментально измеряемые величины без приближений идеального газа и идеального раствора, успешно проверена на системе «метан–этан».
  \item Разработана и апробирована методика расчёта коэффициента термической аккомодации в рамках метода молекулярной динамики. Получены зависимости коэффициента от температуры и размера кластера для наночастиц железа, взаимодействующих с атомами аргона.
  \item Показано, что количество энергии, уносимой от кластера налетающим атомом, зависит от температуры кластера слабее, чем по линейному закону. Установлено, что данная нелинейность обусловлена уменьшением времени взаимодействия атома с кластером при росте температуры.
\end{enumerate}


Последний параграф может включать благодарности.  В заключение автор
выражает благодарность и большую признательность научному руководителю
Иванову~И.\,И. за поддержку, помощь, обсуждение результатов и~научное
руководство. Также автор благодарит Сидорова~А.\,А. и~Петрова~Б.\,Б.
за помощь в~работе с~образцами, Рабиновича~В.\,В. за предоставленные
образцы и~обсуждение результатов, Занудятину~Г.\,Г. и авторов шаблона
*Russian-Phd-LaTeX-Dissertation-Template* за~помощь в оформлении
диссертации. Автор также благодарит много разных людей
и~всех, кто сделал настоящую работу автора возможной.
